\documentclass[11pt,a4paper,twoside]{article}
\usepackage[utf8]{inputenc}
\usepackage[english]{babel}
\usepackage{actuarialangle, actuarialsymbol}


\begin{document}
Describe the setup of the chain ladder method and determining its factors.


Time window of \( n + 1 \) years (\( n \in \mathbb{N} \))
\(\rightarrow n + 1 \) potential years of occurrence
Each claim settled ...
• either in its year of occurrence ...
• or in one of the \( n \) years thereafter
\[
[n] := \{0, 1, 2, \ldots, n\}, \quad n \in \mathbb{N}
\]
Claims occurring in year \( i \) settled in year \( k \) thereafter lead to random payment \( Z_{i,k} \),  
\( i, k \in [n] \)
\[
i + k \leq n \Longrightarrow \text{value } Z_{i,k} \text{ known (cf. triangle)}
\]
\[
i + k > n \Longrightarrow \text{value of } Z_{i,k} \text{ not known yet!}
\]

Goal: Add estimates for the values of \( Z_{i,k},\ i + k > n \), to triangle

Cumulation −→ losses
−→ random variables \( S_{i,k} , i, k ∈ [n] \), \( S_{i,k} := \sum_{l=0} Z_{i,l}\)
As before: \( S_{i,k} \) known if \( i+k \leq n \), not yet known if \( i + k > n \)


Best-known claims reserving method
Idea 1: All accident years behave similarly
Idea 2: \( S_{i,k} \approx \text{factor}_k S_{i,k-1} \)

Consider the quantities
\(
F_{i,k} := \frac{S_{i,k}}{S_{i,k-1}}, \quad i, k \in [n], \quad k \ne 1
\)

Known if \( i + k \le n \)
Note: For \( i + k > n \), we have
\(
S_{i,k} = \underbrace{S_{i,n-i}}_{\text{observable}} \prod_{l = n - i + 1}^{k} F_{i,l}
\)
Estimate the unknown factors!
First step: Estimate \( F_{i,k} \) with
\[
\hat{\varphi}_k^{CL} := \frac{\sum_{j=0}^{n-k} S_{j,k}}{\sum_{j=0}^{n-k} S_{j,k-1}} = \sum_{j=0}^{n-k} \frac{S_{j,k-1}}{\sum_{l=0}^{n-k} S_{l,k-1}} \cdot F_{j,k}
\]
(weighted arithmetic mean, no dependence on \( i \)!)
Second step: Estimate \( \mathbb{E}[S_{i,k}] \) for \( i \in [n] \) and \( k = n - i, \ldots, n \) by
\[
\hat{S}_{i,k}^{CL} := S_{i,n-i} \prod_{l = n - i + 1}^{k} \hat{\varphi}_l^{CL}
\]

Chain-ladder reserve for year of occurence \( i: \hat{R}_{i}^{CL} := \hat{S}_{i,n}^{CL} - S_{i,n} \)
Total chain-ladder reserve: \( \hat{R}^{CL} = \sum_{i=0}^n \hat{R}_{i}^{CL} \)


Describe the setup of the crossing up method and determine its factors.

Reserves only depend on unobservable information about \(S_{i,n}\), \(i \in [n]\)
Alternative: Estimate \(S_{i,n}\) directly with grossing-up method
Set \(\gamma_{i,n-1} = \frac{S_{i,n-1}}{S_{i,n}}\), \(i \in [n]\)
Trivial: 
\(S_{i,n} = \frac{S_{i,n-1}}{\gamma_{i,n-1}}\) 
Compute 
\(\gamma_{i,n-1}\)'s iteratively starting with \(i = 0\): 
\[\hat{\gamma}_{n-i}^{GU} :=
\begin{cases}
1 & i = 0 \\
\frac{\sum_{j=0}^{i-1} S_{j,n-i}}{\sum_{j=0}^{i-1} \hat{S}_{j,n}^{GU}} & i > 0
\end{cases}\] 

\(
\hat{S}_{i,n}^{GU} := \frac{S_{i,n-i}}{\hat{\gamma}_{n-i}^{GU}}, \quad i \in [n]
\)
Grossing-up reserve for year of occurence \( i: \hat{R}_{i}^{GU} := \hat{S}_{i,n}^{GU} - S_{i,n-i} \)
Total grossing-up reserve \( \hat{R}_{i}^{GU} := \sum_{i=0}^n \hat{R}_{i}^{CL} \)

\end{document}

